\documentclass[11pt]{article}
\usepackage[margin=1in]{geometry}
\usepackage{amsmath,amssymb,amsthm,hyperref,enumitem}
\hypersetup{colorlinks=true,linkcolor=blue,urlcolor=blue}
\newtheorem{theorem}{Theorem}
\newtheorem{lemma}{Lemma}
\newcommand{\E}{\mathbb E}
\newcommand{\Pp}{\mathbb P}
\title{Nonsingularity of the periodic partial-sums system\\with arbitrary support minima}
\author{Principia Math In Progress}
\date{22 September 2026}
\begin{document}
\maketitle
\noindent\textbf{MathDB 375984. Full solution of Grigutis--Petkelis Conjecture 1.}

\section*{Task 1: statement and the source conventions}
Let independent integer increments have period \(N\) in distribution,
minimum support points \(m_j\), and
\[
 D=\sum_{j=1}^Nm_j<0,\qquad
 \mu=\E\sum_{j=1}^NX_j<0,\qquad
 \Pp(X_j=m_j)>0.
\]
Indices on phases are cyclic. Write
\(\mathcal M_j=\max(0,X_j,X_j+X_{j+1},\ldots)\) and
\(\pi^{(j)}_i=\Pp(\mathcal M_j=i)\).
The source constructs a system for its initial probabilities by
evaluating generating-function numerators at the nonunit roots
of \(G_{\sum X_j}(z)=1\), and appending a normalization equation
[1, (2.7), pp.~7--9]. It conjectures nonsingularity.

Two source conventions matter. A root of multiplicity \(q\)
gives derivative constraints of orders \(0,\ldots,q-1\)
[1, Note~1]. Probabilities forced to be zero by support bounds
are removed, as in the mixed-sign example [1, pp.~17--18].
The printed block dimensions \(-m_j\) apply directly when
all \(m_j\leq0\); for general integer minima the forced-zero
elimination below specifies the \(-D\) actual unknowns.

\begin{theorem}\label{thm}
Under the assumptions above, the source's coefficient system,
with these support and multiplicity conventions, is nonsingular.
This holds for every period, arbitrary integer support minima,
all multiplicities of interior roots, and roots on the unit circle.
\end{theorem}

\subsection*{1. Reduce all support minima to nonpositive values}
Define deterministic heights
\[
 \ell_j=\max\left(0,m_j,m_j+m_{j+1},\ldots,
                         m_j+\cdots+m_{j+N-1}\right).
\]
Since \(D<0\), an extra full period cannot increase a partial
sum of the minima. Consequently
\[
 \ell_j=\max(0,m_j+\ell_{j+1}),\qquad
 \mathcal M_j\geq\ell_j.
\]
Set
\[
 k_j=\ell_j-m_j-\ell_{j+1}\geq0,\quad
 C=\sum_j k_j=-D>0,\quad
 Y_j=X_j+\ell_{j+1}-\ell_j,\quad
 Z_j=Y_j+k_j=X_j-m_j.
                                                               \tag{1}
\]
Here \(Z_j\geq0\) and \(p_{j,0}:=\Pp(Z_j=0)>0\).
The transformed increments have minima \(-k_j\) and the same
period sum as the original increments. Their means are finite:
each increment is bounded below, and their total mean is finite
and negative. In particular \(\sum_j\E Z_j=C+\mu<C\).

The translated maxima \(L_j=\mathcal M_j-\ell_j\) obey
\[
 L_j\ \buildrel d\over=\ \max(0,Y_j+L_{j+1}).
                                                               \tag{2}
\]
Indeed the original recursion gives
\(L_j=\max(-\ell_j,Y_j+L_{j+1})\) in distribution.
If \(\ell_j=0\), this is (2). If \(\ell_j>0\), then \(k_j=0\)
and \(Y_j\geq0\), so both maxima equal \(Y_j+L_{j+1}\).

In phase \(j\), the source initially uses
\(\pi^{(j+1)}_i\), \(0\leq i<-m_j\). Those with
\(i<\ell_{j+1}\) are zero. If \(\ell_j>0\), none remain.
If \(\ell_j=0\), the remaining indices are precisely
\[
 i=\ell_{j+1}+u,\quad 0\leq u<k_j.
                                                               \tag{3}
\]
Thus exactly \(\sum_jk_j=C=-D\) unknowns remain in all cases.
When every \(m_j\leq0\), all heights vanish and this is
the displayed source matrix without any change of variables.

\subsection*{2. An equivalent description of every matrix row}
Put \(g_j(z)=\E z^{Z_j}\), \(F_j(t)=\Pp(Z_j\leq t)\).
For an arbitrary complex vector \(v=(v_{j,u})_{0\leq u<k_j}\),
define
\[
 Q_j(z)=\sum_{u=0}^{k_j-1}v_{j,u}
                  \sum_{t=u}^{k_j-1}z^tF_j(t-u),
 \qquad Q_j=0\ \text{if }k_j=0.                              \tag{4}
\]
The map from \(v_j\) to the coefficients of \(Q_j\) is triangular
with positive diagonal \(F_j(0)=p_{j,0}\).
Let
\[
 K(z)=\prod_{j=1}^Ng_j(z)-z^C,
 \qquad
 H_j(z)=\sum_{r=0}^{N-1}
 z^{\,C-\sum_{t=0}^r k_{j+t}}
 \left(\prod_{t=0}^{r-1}g_{j+t}(z)\right)Q_{j+r}(z).
                                                               \tag{5}
\]
Empty products equal one. Every exponent in (5) is nonnegative.
The root-count result used in [1], namely [2, Lemma~4] applied
with one aggregate variable \(\sum_jZ_j\) and premium \(C\),
gives \(C\) zeros of \(K\) in the closed unit disk, counted with
multiplicity. One is \(1\). None is zero, since \(K(0)=\prod_jp_{j,0}>0\).

Up to invertible row operations, the homogeneous system is
\[
 H_N^{(t)}(\alpha)=0
 \quad\left(K(\alpha)=0,\ \alpha\ne1,\
                     0\leq t<\operatorname{mult}_K(\alpha)\right),
 \qquad \sum_jQ_j(1)=0.                                    \tag{6}
\]
Here and below boundary roots will be shown to be simple.
To verify the exact source correspondence, let \(B_j^{\rm old}\)
be the Laurent polynomial in [1, Proposition~1].
After substituting (3) into its definition, direct index shifting
gives
\[
 B_j^{\rm old}(z)=z^{\ell_j-k_j}Q_j(z),\qquad
 G_{X_j}(z)=z^{\ell_j-\ell_{j+1}-k_j}g_j(z).                 \tag{7}
\]
The first equality is also true when \(\ell_j>0\), since both
sides vanish. The numerator in the source's phase-\(N\) row,
with block order \(N,1,\ldots,N-1\), is therefore
\[
 B_N^{\rm old}+G_{X_N}B_1^{\rm old}
       +G_{X_N}G_{X_1}B_2^{\rm old}+\cdots
       =z^{\ell_N-C}H_N(z).                                \tag{8}
\]
At a nonzero root, multiplication by \(z^{\ell_N-C}\)
changes a vector of derivatives through order \(q-1\)
by an invertible triangular matrix. The factor \(1-z\)
in the source identity is likewise invertible at \(\alpha\ne1\).
At \(z=1\), (7) gives the same normalization row
\(\sum_jQ_j(1)\), with right-hand side \(-\mu\).
This proves (6), including the confluent rows.

\subsection*{3. Kernel cancellation gives decaying coefficients}
Suppose (6) holds. Telescoping (5) gives
\[
 g_jH_{j+1}-z^{k_j}H_j=KQ_j,\qquad 1\leq j\leq N.
                                                               \tag{9}
\]
At any nonzero root \(\alpha\) of \(K\), all \(g_j(\alpha)\)
are nonzero because their product is \(\alpha^C\).
Equation (9) propagates the prescribed vanishing order from
\(H_N\) to all \(H_j\). Also
\(H_j(1)=\sum_rQ_r(1)=0\).
It follows that \(T_j=H_j/K\) is holomorphic in the open unit disk.

We check continuity on its boundary without imposing exponential
tails or higher moments. Finite first moments imply that \(g_j\)
and \(g_j'\) are continuous on the closed disk, by uniform absolute
convergence of their series. The same holds for \(K,H_j\) and
their first derivatives.

If \(|\alpha|=1\) and \(K(\alpha)=0\), let \(Z=\sum_jZ_j\).
Equality in \(|\E\alpha^Z|\leq1\), together with \(\Pp(Z=0)>0\),
forces \(\alpha^r=1\) on the support of \(Z\), and
\(\alpha^C=1\). Therefore
\[
 K'(\alpha)=\alpha^{-1}(\E Z-C)=\alpha^{-1}\mu\ne0.          \tag{10}
\]
All boundary roots are simple. If a function \(f\) and \(f'\)
are continuous on the closed disk and \(f(\alpha)=0\),
integration along the segment from \(\alpha\) to \(z\) shows
\[
 \frac{f(z)}{z-\alpha}\longrightarrow f'(\alpha)
 \quad\text{as }z\to\alpha\text{ within the disk}.
 \]
Apply this to \(H_j,K\); (10) proves continuous extension
of \(T_j\) at every boundary root.

Consequently the Taylor coefficients
\(h_j(u)=[z^u]T_j(z)\) tend to zero. For completeness, if
\(T\) is holomorphic in the disk and continuous on its closure,
Cauchy's estimate bounds its coefficients by
\(M=\max|T|\). For \(0<\rho<1\), apply that estimate to
\(T(z)-T(\rho z)\) to get
\[
 |h(u)|\leq M\rho^u+
      \max_{|z|\leq1}|T(z)-T(\rho z)|.
 \]
First let \(u\to\infty\), then \(\rho\uparrow1\).
Uniform continuity proves the assertion.

\subsection*{4. A cyclic maximum principle finishes the proof}
Dividing (9) by \(K\) gives
\[
 g_jT_{j+1}-z^{k_j}T_j=Q_j.
\]
Since \(\deg Q_j<k_j\) (with \(Q_j=0\) when \(k_j=0\)),
coefficient comparison yields, for every \(u\geq0\),
\[
 h_j(u)=\sum_{r=0}^{u+k_j}p_{j,r}\,
                        h_{j+1}(u+k_j-r).                 \tag{11}
\]
Here \(p_{j,r}=\Pp(Z_j=r)\).
If these finitely many decaying sequences were not all zero,
their joint supremum in modulus would be a positive attained
maximum \(M_*\), say \(|h_j(u)|=M_*\).
The right side of (11) has modulus at most \(M_*\), because
the weights are nonnegative and sum to at most one.
Since \(p_{j,0}>0\), equality forces
\(|h_{j+1}(u+k_j)|=M_*\).
Repeating this implication for a complete period gives
\(|h_j(u+C)|=M_*\), then \(|h_j(u+qC)|=M_*\) for every \(q\).
As \(C>0\), this contradicts decay.

Thus all \(T_j\), all \(Q_j\), and, by the triangular
invertibility in (4), all \(v_{j,u}\) vanish. The square
system of \(C\) rows has trivial kernel. Equations (7)--(8)
transfer this conclusion to the original system, proving
Theorem~\ref{thm}.

\section*{Scope and attribution}
No condition that the source's boundary terms \(b_k(x)\) vanish
was used. The proof covers unequal negative minima, zero minima,
positive minima after the exact support reduction, and every
root multiplicity allowed by Note~1. It does not treat duplicated
evaluation rows as the intended repeated-root system.

The root count is the established result [2, Lemma~4].
The analytic cancellation and maximum principle extend the
earlier campaign proof for MathDB 357004 [3], whose theorem
assumed a common positive premium and simple nontrivial roots.
That earlier result is credited as the method's source.
The additional support-height reduction, variable and zero shifts,
and all confluent cases are proved here.

\section*{Tasks 2--5: proposed mathematical follow-ups}
\begin{enumerate}[label=\textbf{Task \arabic*:},start=2,leftmargin=*]
\item Obtain quantitative lower bounds for the smallest singular
value of the normalized confluent matrix in terms of the drift,
minimum masses, root separation, and moment bounds.
\item Extend the uniqueness theorem to a finite-state Markov
environment and identify the appropriate drift and irreducibility
conditions for a matrix-valued kernel.
\item Determine the rank and compatibility conditions at zero
period drift, including the degenerate case of an identically
zero period sum.
\item Construct a stable algorithm that avoids explicit root
locations while computing the same initial probabilities, with
rigorous error bounds for unbounded increment distributions.
\end{enumerate}

\begin{thebibliography}{9}
\bibitem{gp}
A.~Grigutis and J.~Petkelis, \emph{The limit law of the maximum
of discrete partial-sums distribution II},
\href{https://arxiv.org/abs/2607.11275}{arXiv:2607.11275v1}
(13 July 2026), Conjecture~1, Note~1, and Example~1.
\bibitem{gg}
S.~Gerv\'e and A.~Grigutis, \emph{Distribution of shifted discrete
random walk generated by distinct random variables and applications
in ruin theory}, Modern Stochastics: Theory and Applications
\textbf{11}(3) (2024), 323--357, Lemma~4, pp.~338--339.
\href{https://doi.org/10.15559/24-VMSTA249}{doi:10.15559/24-VMSTA249}.
\bibitem{earlier}
Principia Math In Progress, \emph{Nonsingularity of the periodic
random-walk system}, MathDB 357004, 18 September 2026.
\href{https://github.com/antoshashakov/Principia-Math-In-Progress/blob/mathdb-results/mathdb-open-problems/problems/357004/solution.tex}{Earlier campaign proof}.
\end{thebibliography}
\end{document}
